\lecture{8}
\title{Confidence Intervals}

\begin{document}

\subsection{Asymptotic Normality}

\textbf{Definition.} An estimator $\hat{\theta}_n$ of $\theta$ is \textbf{asymptotically normal}
if there is some \textbf{asymptotic variance} $\sigma^2 > 0$ such that:
\[ \sqrt{n}(\hat{\theta}_n - \theta) \rightsquigarrow \mathcal{N}(0, \sigma^2). \]
Informally, this is roughly equivalent to
$\hat{\theta}_n \approx \theta + \frac{\sigma}{\sqrt{n}}Z$, where $Z \sim \mathcal{N}(0, 1)$.

\begin{problem}
    Given model $\{\mathrm{Ber}(p) \mid p \in (0, 1)\}$
    and estimator $\hat{\theta}_n := \bar{X}_n$ of $\theta := p$,
    determine its asymptotic variance.
\end{problem}

\begin{solution}
    Since $\mathbb{E}_p[X_i] = p = \theta$ and
    $\mathbb{V}_p[X_i] = p(1 - p)$, the CLT says that
    the asymptotic variance is $p(1 - p)$.
\end{solution}

\begin{problem}
    Given model $\{\mathrm{Ber}(p) \mid p \in (0, 1)\}$
    and estimator $\hat{\theta}_n := (\bar{X}_n)^2$ of $\theta := p^2$,
    determine its asymptotic variance.
\end{problem}

\begin{solution}
    Invoke the delta method with $g(x) = x^2$;
    the asymptotic variance gets scaled by
    $g'(p)^2 = 4p^2$, so it's $4p^3(1 - p)$.
\end{solution}

\begin{remark}
    If $\hat{\theta}_n$ is asymptotically normal
    with asymptotic variance $\sigma^2$,
    that tells us that
    $\hat{\theta}_n \approx \mathcal{N}(\theta, \frac{\sigma^2}{n})$.

    However, that does \emph{not} imply that
    the variance of $\hat{\theta}_n$ equals---or even
    \emph{converges to}---the variance of
    $\mathcal{N}(\theta, \frac{\sigma^2}{n})$.
    Convergence in distribution does not imply
    convergence of moments.

    Here's an explicit counterexample:
    consider an estimator $\hat{\theta}_n$ that happens to satisfy
    $\hat{\theta}_n := \theta + \frac{Y_n}{\sqrt{n}}$,
    with the random variable $Y_n$ defined as below.
    \[ Z \sim \mathcal{N}(0, \sigma^2) ~~~ \parallel ~~~
    Y_n := \begin{cases}
    Z & \text{ with probability } 1 - \frac{1}{n}. \\
    n & \text{ with probability } \frac{1}{n}.
    \end{cases} \]
    Then we have a convergence in distribution
    $Y_n \rightsquigarrow \mathcal{N}(0, \sigma^2)$, implying that
    $\sqrt{n}(\hat{\theta}_n - \theta) \rightsquigarrow
    \mathcal{N}(0, \sigma^2)$.
    In other words, the estimator $\hat{\theta}_n$
    is asymptotically normal with asymptotic variance $\sigma^2$.
    However, we may compute
    $\mathbb{V}[\hat{\theta}_n]
    = \frac{1}{n} \cdot \mathbb{V}[Y_n] \approx 1$ as
    $n \to \infty$, which
    does not match the asymptotic variance $\sigma^2/n \to 0$.

    That said, in most well-behaved cases (read: under regularity conditions),
    it will generally be true that
    $\mathbb{V}[\hat{\theta}_n] \approx \frac{\sigma^2}{n}$.
    So it is still generally fine to intuit
    $\sigma^2$ and $n \cdot \mathbb{V}[\hat{\theta}_n]$
    as the same things.
\end{remark}

\subsection{Confidence Intervals}

\textbf{Definition.} We say a $(1 - \alpha)$ \textbf{confidence interval (CI)}
for a parameter $\theta$ is a \emph{random} interval of the form
$C_n = (A_n, B_n)$ such that, for any fixed value of $\theta$,
we have $\Pr_{\theta}(\theta \in (A_n, B_n)) \geq 1 - \alpha$.
We say $1 - \alpha$ is the \textbf{coverage} of the CI.

\begin{remark}
    For any fixed $\theta$, the endpoints $A_n$ and $B_n$ are themselves
    random variables whose distributions (usually) depend on $\theta$.

    If it turns out $(6, 7)$ is a $95\%$ confidence interval for $\theta$,
    that does not mean $\Pr_{\theta}(\theta \in (6, 7)) \geq 0.95$.
    That statement doesn't even make sense---either $\theta \in (6, 7)$
    or $\theta \not \in (6, 7)$, and there (usually) is no distribution on
    the possible values of $\theta$.
\end{remark}

\textbf{Definition.} A sequence of random intervals
$\{C_n\}_{n = 1}^{\infty}$ has \textbf{asymptotic coverage} $1 - \alpha$
if $\liminf_{n \to \infty} \Pr_{\theta}(\theta \in C_n) \geq 1 - \alpha$.

\begin{theorem}{Constructing CIs}
    Suppose $\hat{\theta}_n$ is an asymptotically normal estimator of $\theta$
    with asymptotic variance $\sigma^2$.  Then:
    \[ C_n := \left( \hat{\theta}_n - z_{\alpha / 2} \frac{\sigma}{\sqrt{n}},
    \hat{\theta}_n + z_{\alpha / 2} \frac{\sigma}{\sqrt{n}} \right) \]
    has asymptotic coverage $(1 - \alpha)$, where $z_{\alpha / 2}$
    is such that $\Phi(z_{\alpha / 2}) - \Phi(-z_{\alpha / 2}) = 1 - \alpha$.
\end{theorem}

\begin{proof}
    Just note that $\theta \in C_n$ is equivalent to
    $|\hat{\theta}_n - \theta| < z_{\alpha / 2} \frac{\sigma}{\sqrt{n}}$,
    but also $\hat{\theta}_n - \theta \approx \mathcal{N}(0, \frac{\sigma^2}{n})$.
    ~ $\blacksquare$
\end{proof}

\begin{problem}
    Given (i.i.d.) samples $X_1, \dots, X_n \sim \mathrm{Ber}(p)$,
    construct confidence intervals for $p$ with asymptotic coverage $1 - \alpha$.
\end{problem}

\begin{solution}
    The challenge is that $\sigma = \sqrt{p(1 - p)}$ depends on $p$,
    but we don't know $p$.
    
    The fix is to conservatively bound $\sigma \leq \frac{1}{2}$, which yields
    $C_n = \bar{X}_n \pm \frac{z_{\alpha / 2}}{2\sqrt{n}}$.
\end{solution}

\end{document}